\section{5. Comparison with existing methods}\label{comparison-with-existing-methods}

Table 4 places FLEX-DBS alongside conventional tethered stimulation and published wireless rodent stimulators that are small enough for mice or are architecturally instructive. The columns are the properties that determine whether a system can run a given experiment: species and mass, channels, output type, supply or compliance, parameter range, how parameters are set, power and operating duration, and whether the hardware and software are released.

\begin{landscape}
\textbf{Table 4.} FLEX-DBS in the context of tethered stimulation and published wireless rodent stimulators. Values for published systems are as reported by the authors; blank cells are not reported in the source. CC, constant current; CV, constant voltage.

{\scriptsize
\setlength{\tabcolsep}{2pt}
\begin{longtable}{@{}>{\raggedright\arraybackslash{}}p{0.09\linewidth}>{\raggedright\arraybackslash{}}p{0.07\linewidth}>{\raggedright\arraybackslash{}}p{0.08\linewidth}>{\raggedright\arraybackslash{}}p{0.08\linewidth}>{\raggedright\arraybackslash{}}p{0.05\linewidth}>{\raggedright\arraybackslash{}}p{0.10\linewidth}>{\raggedright\arraybackslash{}}p{0.14\linewidth}>{\raggedright\arraybackslash{}}p{0.12\linewidth}>{\raggedright\arraybackslash{}}p{0.12\linewidth}>{\raggedright\arraybackslash{}}p{0.09\linewidth}@{}}
\toprule
System & Species tested & Mass with battery & Channels & Output & Supply / compliance & Amplitude, pulse width, frequency; waveform & Parameter control & Current draw; operating duration & Open release \\
\midrule
\endhead
Tethered stimulator (reference method) & Any & n/a (tether) & Instrument-dependent & CC or CV & Instrument-dependent & Instrument-dependent & Wired, real time & Mains powered; limited by the tether & n/a \\
\textbf{FLEX-DBS (this work)} & Mouse & & 2; shared amplitude and timing & CC & ±5 V rails; $\sim$4.9 V swing & 0--625 µA ideal (compliance-limited in vivo); $\sim$117--600 µs; 83.3--166.7 Hz delivered; charge-balanced biphasic, 1:5 recovery; continuous or burst & Real-time BLE from an iOS app; autonomous after disconnect & 3.15--3.42 mA stimulating; $\sim$3 d per cell pair & Schematic, firmware, app, protocol (S1--S4) \\
t-IPG, \citet{paulat2026} & Mouse (s-IPG also in rats and mice) & 0.9 g (s-IPG 1.6 g) & 2, time-delayed & CC & 3 V lithium cell, direct & 1--100 µA, 50--700 µs, 2--257 Hz; cathodic phase with long low-amplitude anodic phase; theta- and beta-burst modes & Predefined paradigms via near-field transponder & 21.3 µA maximum; 49 d voltage-compliant & Partial (GitHub) \\
STELLA, \citet{plocksties2021} & Rat (proposed for mice $>$20 g) & 4.0 g encapsulated (CR1225); 2.6 g (CR1216) & 2, synchronized & CC & Boost 3.1--3.7 V; 2.2--2.8 V at electrode & 42--100 µA, 1 µs--100 \% duty, 10--5000 Hz; monophasic, passive charge balance & Firmware set before implantation; Hall sensor steps through presets & 7.6--8.3 µA; $\ge$12 wk in vitro, 42 d in vivo & Yes, open license (GitHub) \\
\citet{grotemeyer2024} & Mouse & 2 g & 1 & CC & 3 V (two 1.5 V cells) & 10--500 µA, 60--500 µs, 10--200 Hz; monophasic, passive discharge & Cable programmer, reprogrammable after implantation; reed switch on/off & $\sim$20 µA; 52 d in vitro & Analysis code only; commercialized \\
\citet{fleischer2020} & Mouse & 2.8 g & 1 & CC & 3 V (two V364 cells) & 0--500 µA (text: up to 300 µA), 60--500 µs, 10--200 Hz (text: 10--500 Hz); monophasic, passive discharge & Programmed before mounting; switch on/off & $\sim$20 µA; $\sim$30 d in vitro & \\
\citet{dehaas2012} & Mouse & & 2 (bilateral) & & & Biphasic & & & \\
\citet{alpaugh2019} & Mouse & 4.7 g & 1 & CC & 3 V (two 393 cells); 2.1 V output headroom & Fixed 149 µA, $\sim$103 µs, $\sim$130 Hz; balanced biphasic & Preprogrammed via external circuit; no radio & $\sim$2 wk per cell set, $\ge$30 d with weekly exchange & Component list only \\
\citet{kouzani2013} & Rat (designed for mice) & 3.27 g (head-mount) & 1 & CC & 3.1 V (two silver-oxide cells) & Fixed 200 µA, 90 µs, 130 Hz; monophasic, capacitor-coupled & Preprogrammed; changes require reflashing or a resistor change & $\sim$870 µA; 12.3 d on bench, 7 d in vivo & Schematic and pseudocode \\
\citet{pinnell2018} & Rat & 2.8 g with battery and head cap & 2 & CC & 12.29 V regulated & 20 µA--2 mA, 10 µs--100 \% duty, 0.1--5000 Hz; monophasic or biphasic; dual-frequency burst & Wired programmer; amplitude by potentiometer; magnet on/off & 35 µA in low-power state; $\sim$30 h stimulating & Yes (Figshare) \\
\citet{burton2021} & Rat & 0.046 g, battery-free & 1 bipolar electrode & CV & 1.5--5.5 V (regulated) & 1.5--5.5 V, $\ge$20 µs pulses; biphasic & Real time via modulation of RF power & 9.35 mW while stimulating; $>$36 d in vivo with RF power & \\
\citet{shen2026} & Rat & 12.5 g (7.5 g battery) & & CV & 3.7 V Li-ion & 2.3--3 V, 40--340 µs, 80--150 Hz; biphasic & Real-time BLE (MATLAB interface) & $\ge$300 h on 400 mAh & \\
\bottomrule
\end{longtable}
}
\end{landscape}

Three patterns emerge. First, the mouse-scale systems with runtimes of weeks \citep{paulat2026,plocksties2021,grotemeyer2024,fleischer2020} run from supplies of about 3 V, either directly from the cells or through a small boost that leaves 2.2--2.8 V at the electrode, and draw roughly 8--22 µA. Their parameters are set before implantation or through a near-field or cable programmer, and a magnet or reed switch turns stimulation on and off or steps through presets. Their average currents are two orders of magnitude below the 3.15--3.42 mA of FLEX-DBS, and that gap reflects the cost of compliance rails, precision analog circuitry and a live radio rather than of stimulation. The low supply also bounds the load these systems can drive: \citet{alpaugh2019} note that their 2.1 V of output headroom delivers the full 149 µA only into loads below about 20 kΩ, well below the 37--80 kΩ measured in vivo in Section 3. Second, the systems that offer real-time wireless control are rat-scale or externally powered and regulate voltage rather than current: \citet{burton2021} harvest power from a radio-frequency field in the arena, and \citet{shen2026} use BLE with a 7.5 g battery. \citet{pinnell2018} combine constant current, two channels and 12.29 V compliance in a 2.8 g package, but set parameters through a wired programmer and potentiometers and run for about 30 h. Third, no system in Table 4 combines bilateral constant-current output, compliance rails of ±5 V, continuously programmable amplitude, pulse width and frequency, burst mode, and real-time phone control. That combination is what FLEX-DBS contributes, and it is bought with runtime measured in days rather than weeks.

The comparison is not a claim of universal superiority. For a fully implanted study with fixed parameters and no routine access to the animal, an architecture such as that of \citet{paulat2026} or \citet{plocksties2021} is suitable. FLEX-DBS is designed for experiments in which the investigator must set, change and verify parameters at the cage, deliver known current into an uncertain load, and stimulate continuously from hours through multiple days. Section 6 tests whether it behaves like the tethered reference method it replaces.
